\documentclass{article}
\usepackage[T1]{fontenc}
\usepackage{lmodern}
\usepackage{graphicx}
\usepackage{amsmath,amssymb}
\usepackage[a4paper,margin=2cm]{geometry}
\usepackage[absolute,overlay]{textpos}
\setlength{\TPHorizModule}{1mm}
\setlength{\TPVertModule}{1mm}
\usepackage[indonesian]{babel}
\usepackage{hyperref}
\setlength{\emergencystretch}{2em}
% unit: Halaman 9

\begin{document}

% === Halaman 9 (llm) ===

\begin{itemize}
    \item Ini sering terjadi jika sebagian isi pesan bisa ditebak (misalnya \textit{header file}, protokol standar, salam pembuka email).
    \item Beberapa pesan yang formatnya terstruktur membuka peluang untuk menerka plainteks dari cipherteks yang bersesuaian.
\end{itemize}

Contoh:

\begin{itemize}
    \item \textit{From} dan \textit{To} di dalam e-mail,
    \item "Dengan hormat", \textit{wassalam}, pada surat resmi.
    \item \textit{\#include}, \textit{program}, di dalam \textit{source code}
\end{itemize}

\vspace{10mm}

\begin{itemize}
    \item Dengan menggunakan pasangan plainteks dan cipherteks, kriptanalis dapat menemukan kunci enkripsi (lihat pembahasan \textit{cipher} klasik \textit{affine cipher} dan \textit{hill cipher})
\end{itemize}

\end{document}
